\section{Engineering findings}
\label{app:engineering}

These findings concern the system built around the benchmark rather than the benchmark itself. They
are moved here from the main text so that it stays the benchmark audit; only the opening sentence of
each is adapted to its new place, except the served-constants finding, whose figures and the claim
they carried were removed because no run output was kept for them.

\subsection*{A fix that did nothing in its own scenario}

The defect of Section~\ref{sec:discussion}'s ``a test that cannot fail'' appeared in a fix rather
than a test. A concurrency fix in the
staff-identity service (a refresh after a row-locking query, so that a stale copy of an account could
not be written back) passed its test and was recorded as fixed, with a mutation killed. The test
created the conflict in a way that never advanced the row's version, so it exercised only the branch
where the fix was not needed; in the case the fix was written for, the locking query itself raised an
error before the refresh ran. A second defence made every assertion hold, and the killed mutation
removed both defences at once. A mutation removing only the refresh survived, and was first filed as
``defence in depth''; the case was found when that explanation was required to be shown with a test
that fails without the fix, and the test failed with the fix in place. Silent vacuous tests are bad; a
record that reads ``fixed, tested, mutation killed'' for a fix that did nothing in its own scenario is
worse. Counting killed mutations measures the tests; only investigating the survivors measures the
claims, and an ``equivalent'' or ``defence in depth'' label on a survivor needs the same evidence as
any other claim. That last rule is written, not enforced.

\subsection*{Served features the evaluation did not see}
 The scoring service as built (milestone
M5) serves \nSkewConstantFeatures{} of the trained features as constants, because nothing deployed yet
writes the state they read: the counterparty's labelled-fraud count, the map cell's lagged fraud
rate and the days since a SIM swap; the synthetic-identity score is computed from fewer of its terms
(ADR 0034). ADR 0034 and the scoring milestone's re-review record figures for what this does to
the evaluated model on the test period, but no run output was kept for it, so it is not reported
here; whether the gate's metrics would have registered the change is therefore not established. The
serving path carries its own acceptance test, and M10 is required to re-measure the skew and find
none (PB-72).

\subsection*{The verification effort was blind in one direction}
 The scoring client was reviewed
\nFakeReviews{} times (a principal review, a re-review of its fixes, a verification pass and an
adversarial pass over the promotion path, each by an agent), and the reviewers found \nFakeDefects{} defects in the
model-promotion path between them, while the suite passed locally with its real-registry test skipped,
since nothing that runs without Docker touches a real server (\texttt{docs/\allowbreak parallel/\allowbreak
M5\_updates.md}). None found that the client misread how the model registry reports an unset alias,
because the tests of that path that ran without Docker used an in-process fake that encoded the same
misreading. The test that starts a real registry container caught it in continuous integration: in that run it was the suite's
\nFakeCiFailed{} failure while \nFakeCiPassed{} tests passed
(\texttt{evidence/\allowbreak mlflow\_\allowbreak alias\_\allowbreak ci\_\allowbreak 90078b2.txt}, the
run's whole log). Review, mutation testing and
adversarial probing all operate inside the world the test fixtures define; contact with the real
dependency is the only check outside it. The verification reported for this system therefore does
not establish correctness against the external services it integrates with, beyond the tests that
start real ones, and the scoring milestone's re-review records that some of those had not run on any
machine visible to it.
